\documentclass[11pt,a4paper]{article}
\usepackage[margin=2.5cm]{geometry}
\usepackage[T1]{fontenc}
\usepackage[utf8]{inputenc}
\usepackage{booktabs}
\setlength{\parskip}{6pt}
\setlength{\parindent}{0pt}

%% ------------------------------------------------------------------
%% Problem Solving Protocol -- PROSE IS NOT WRITTEN HERE.
%% BiB does not accept papers drafted by LLMs. The authors write the
%% text from the claim register; every section below is a placeholder
%% carrying its source and word budget. RE-NUMBER the bibliography by
%% first appearance once prose is placed (BiB numbering rule).
%% ------------------------------------------------------------------

\title{Microbiome-CDM: An LLM-driven common data model and knowledge
extraction pipeline for human microbiome literature}
\author{<<authors>>}
\date{}

\begin{document}
\maketitle

\section*{Abstract}
PLACEHOLDER (plan abstract row, \textasciitilde200 words): the problem, the protocol, three headline results.
% authors write; drafts anchor: paper/publication/manuscript_draft/

\paragraph{Keywords.} microbiome; common data model; named entity recognition; large language models; text mining; data provenance % up to 6

\section{Introduction: the problem}
PLACEHOLDER (~550 words per /home/jkalra\_hpc/Microbiome-CDM/paper/publication/manuscript\_bib/BIB\_MANUSCRIPT\_PLAN.md, section map 2026-11-03) --- draft 00\_introduction.md; gap vs BugSigDB, MiDRED, DiMB-RE, MINERVA, Higashi et al.; claim IDs R1-R6.\\
BiB policy: prose written by the authors only.

\section{Protocol overview}
PLACEHOLDER (~350 words per /home/jkalra\_hpc/Microbiome-CDM/paper/publication/manuscript\_bib/BIB\_MANUSCRIPT\_PLAN.md, section map 2026-11-03) --- draft 01\_methods...; Figure 1; C01-C08; claim IDs C01-C08.\\
BiB policy: prose written by the authors only.

\section{Stage 1: Corpus and sectioning}
PLACEHOLDER (~250 words per /home/jkalra\_hpc/Microbiome-CDM/paper/publication/manuscript\_bib/BIB\_MANUSCRIPT\_PLAN.md, section map 2026-11-03) --- draft 01\_methods... Corpus.\\
BiB policy: prose written by the authors only.

\section{Stage 2: Entity extraction with multiple programs}
PLACEHOLDER (~550 words per /home/jkalra\_hpc/Microbiome-CDM/paper/publication/manuscript\_bib/BIB\_MANUSCRIPT\_PLAN.md, section map 2026-11-03) --- draft 01\_methods... Entity mention extraction + NER benchmark; Table 1; claim IDs H03-H07, L01-L04.\\
BiB policy: prose written by the authors only.

\section{Stage 3: Normalisation and vocabulary mapping}
PLACEHOLDER (~400 words per /home/jkalra\_hpc/Microbiome-CDM/paper/publication/manuscript\_bib/BIB\_MANUSCRIPT\_PLAN.md, section map 2026-11-03) --- draft 01\_methods... Canonical terms; guarded mapping, two-model adjudication; claim IDs T, P.\\
BiB policy: prose written by the authors only.

\section{Stage 4: Catalogue and CDM}
PLACEHOLDER (~450 words per /home/jkalra\_hpc/Microbiome-CDM/paper/publication/manuscript\_bib/BIB\_MANUSCRIPT\_PLAN.md, section map 2026-11-03) --- draft 01\_methods...; Figure 2; claim IDs K, M.\\
BiB policy: prose written by the authors only.

\section{Stage 5: Evaluation without human annotation}
PLACEHOLDER (~700 words per /home/jkalra\_hpc/Microbiome-CDM/paper/publication/manuscript\_bib/BIB\_MANUSCRIPT\_PLAN.md, section map 2026-11-03) --- draft 01\_methods... Computational validation + 05\_supplement... Method; Figure 3; rule: concordance, not accuracy; claim IDs N01-N11.\\
BiB policy: prose written by the authors only.

\section{Results: what each program contributed}
PLACEHOLDER (~900 words per /home/jkalra\_hpc/Microbiome-CDM/paper/publication/manuscript\_bib/BIB\_MANUSCRIPT\_PLAN.md, section map 2026-11-03) --- draft 02\_results.md 3-7; 04/05 supplements; Figure 4, Table 2; claim IDs L05-L12, E, J, N02-N08.\\
BiB policy: prose written by the authors only.

\section{Recommendations}
PLACEHOLDER (~300 words per /home/jkalra\_hpc/Microbiome-CDM/paper/publication/manuscript\_bib/BIB\_MANUSCRIPT\_PLAN.md, section map 2026-11-03) --- new interpretive text; Table 3.\\
BiB policy: prose written by the authors only.

\section{Limitations}
PLACEHOLDER (~300 words per /home/jkalra\_hpc/Microbiome-CDM/paper/publication/manuscript\_bib/BIB\_MANUSCRIPT\_PLAN.md, section map 2026-11-03) --- draft 03\_discussion.md Limitations + Structure is not semantics.\\
BiB policy: prose written by the authors only.

\section*{Endmatter (placeholders)}
\begin{itemize}
\item Key Points (3-5 sentences, plan §5)
\item Data availability (deposited bundles)
\item Funding
\item Conflict of interest
\item AI-use statement (required by OUP policy)
\item Reference list (below; BiB style)
\end{itemize}

\section*{Figures (slots)}
Slot for \texttt{fig1\_protocol.pdf} --- Protocol flowchart: 5 stages, programs per stage, counts on arrows\\
Slot for \texttt{fig\_overview.pdf} --- CDM architecture and provenance spine (Figure 2 slot; plan FIGURE\_PLAN F3)\\
Slot for \texttt{figE\_benchmark\_transfer.pdf} --- Evaluation layers (Figure 3 slot)\\
Slot for \texttt{fig4\_results.pdf} --- Per-program results (Figure 4 slot)

\section*{Tables (slots)}
Table 1: programs and their roles (from tables/table1\_programs.csv -- plan Table 1);\\Table 2: results by program (plan Table 2);\\Table 3: recommendations / decision guide (plan Table 3).

\section*{Verification accounting (verified data for the authors; regenerate from the merged evidence file)}

\begin{table}[h]
\centering
\begin{tabular}{lrr}
\toprule
Evidence-anchored verdict (2026-11-07, \emph{DeepSeek-V4.1-Flash}) & $n$ & share \\
\midrule
SUPPORTED & 14{,}363 & 55.7\% \\
PARTIAL & 5{,}780 & 22.4\% \\
NOT\_SUPPORTED & 4{,}618 & 17.9\% \\
CONFLICTS & 860 & 3.3\% \\
NO\_EVIDENCE (quote judged unusable) & 93 & 0.4\% \\
unjudged (69 no-evidence rows; 5 unparsed after retries) & 74 & 0.3\% \\
\midrule
total canonical triples & 25{,}788 & 100\% \\
\bottomrule
\end{tabular}
\caption{Model judgements, not accuracy and not human validation.
Source: \texttt{outputs/results/evidence\_judge\_all\_2026-11-07/REPORT.md} +
\texttt{merged\_triples.csv.gz} (string-level judges of 2026-11-03 and 2026-11-07 are context
columns there).}
\end{table}
\begin{thebibliography}{55}
%
%% GENERATED 2026-11-07 from
%% ../submission_readiness/references_source/bibliography.json (Crossref/PubMed/arXiv/HF/GitHub API, 56 records, 0 errors).
%% Provisional numbering = pipeline-phase order; RENUMBER by first appearance
%% before submission and cite [n] (before punctuation) in the prose.
%%
\bibitem{abdill2025}
Abdill R, Graham S, Rubinetti V et al. Integration of 168,000 samples reveals global patterns of the human gut microbiome. \emph{Cell} 2025;188:1100–1118.e17.
\bibitem{goncalves2019}
Gonçalves R, Musen M. The variable quality of metadata about biological samples used in biomedical experiments. \emph{Sci Data} 2019;6:190021.
\bibitem{geistlinger2023}
Geistlinger L, Mirzayi C, Zohra F et al. BugSigDB captures patterns of differential abundance across a broad range of host-associated microbial signatures. \emph{Nat Biotechnol} 2023;42:790–802.
\bibitem{langarica2025minerva}
Langarica S, Kim Y, Alkhadrawi A et al. MINERVA—microbiome network research and visualization atlas: a scalable knowledge graph for mapping microbiome-disease associations. \emph{Brief Bioinform} 2025;26:bbaf472.
\bibitem{chirumbolo2025minerva}
Chirumbolo S. Methodological and statistical concerns in MINERVA microbiome-disease knowledge graph. \emph{Brief Bioinform} 2025;26:bbaf673.
\bibitem{langarica2025response}
Langarica S, Kim Y, Alkhadrawi A et al. Response to ‘Methodological and statistical concerns in MINERVA microbiome–disease knowledge graph’ by Salvatore Chirumbolo. \emph{Brief Bioinform} 2025;26:bbaf671.
\bibitem{hogan2024midred}
Hogan W, Bartko A, Shang J et al. MiDRED: An Annotated Corpus for Microbiome Knowledge Base Construction. In: \emph{Proceedings of the 23rd Workshop on Biomedical Natural Language Processing}, 2024, 398–408.
\bibitem{hong2025dimbre}
Hong G, Hindle V, Veasley N et al. DiMB-RE: mining the scientific literature for diet-microbiome associations. \emph{J Am Med Inform Assoc} 2025;32:998–1006.
\bibitem{higashi2024harmonization}
Higashi K, Nakagawa Z, Yamada T et al. Automated Harmonization and Large-Scale Integration of Heterogeneous Biomedical Sample Metadata Using Large Language Models. \emph{medRxiv} 2024;2024.10.26.620145 (preprint).
\bibitem{chen2025llmbench}
Chen Q, Hu Y, Peng X et al. Benchmarking large language models for biomedical natural language processing applications and recommendations. \emph{Nat Commun} 2025;16:3280.
\bibitem{li2016bc5cdr}
Li J, Sun Y, Johnson R et al. BioCreative V CDR task corpus: a resource for chemical disease relation extraction. \emph{Database (Oxford)} 2016;2016:baw068.
\bibitem{luo2022biored}
Luo L, Lai P, Wei C et al. BioRED: a rich biomedical relation extraction dataset. \emph{Brief Bioinform} 2022;23:bbac282.
\bibitem{wei2024pubtator3}
Wei C, Allot A, Lai P et al. PubTator 3.0: an AI-powered literature resource for unlocking biomedical knowledge. \emph{Nucleic Acids Res} 2024;52:W540–W546.
\bibitem{sanger2024hunflair2}
Sänger M, Garda S, Wang X et al. HunFlair2 in a cross-corpus evaluation of biomedical named entity recognition and normalization tools. \emph{Bioinformatics} 2024;40:btae564.
\bibitem{akbik2019flair}
Akbik A, Bergmann T, Blythe D et al. FLAIR: An Easy-to-Use Framework for State-of-the-Art NLP. In: \emph{Proceedings of the 2019 Conference of the North American Chapter of the Association for Computational Linguistics (Demonstrations)}, 2019, 54–59.
\bibitem{neumann2019scispacy}
Neumann M, King D, Beltagy I et al. ScispaCy: Fast and Robust Models for Biomedical Natural Language Processing. In: \emph{Proceedings of the 18th BioNLP Workshop and Shared Task}, 2019, 319–327.
\bibitem{cohen2017craft}
Cohen K, Verspoor K, Fort K et al. The Colorado Richly Annotated Full Text (CRAFT) Corpus: Multi-Model Annotation in the Biomedical Domain. In: \emph{Handbook of Linguistic Annotation}. Dordrecht: Springer, 2017, 1379–1394.
\bibitem{kim2004jnlpba}
Kim J, Ohta T, Tsuruoka Y et al. Introduction to the bio-entity recognition task at JNLPBA. In: \emph{Proceedings of the International Joint Workshop on Natural Language Processing in Biomedicine and its Applications - JNLPBA '04}, 2004, 70.
\bibitem{pyysalo2015bionlp13}
Pyysalo S, Ohta T, Rak R et al. Overview of the Cancer Genetics and Pathway Curation tasks of BioNLP Shared Task 2013. \emph{BMC Bioinformatics} 2015;16:S2.
\bibitem{lee2019biobert}
Lee J, Yoon W, Kim S et al. BioBERT: a pre-trained biomedical language representation model for biomedical text mining. \emph{Bioinformatics} 2019;36:1234–1240.
\bibitem{beltagy2019scibert}
Beltagy I, Lo K, Cohan A. SciBERT: A Pretrained Language Model for Scientific Text. In: \emph{Proceedings of the 2019 Conference on Empirical Methods in Natural Language Processing and the 9th International Joint Conference on Natural Language Processing (EMNLP-IJCNLP)}, 2019, 3613–3618.
\bibitem{zaratiana2024gliner}
Zaratiana U, Tomeh N, Holat P et al. GLiNER: Generalist Model for Named Entity Recognition using Bidirectional Transformer. In: \emph{Proceedings of the 2024 Conference of the North American Chapter of the Association for Computational Linguistics: Human Language Technologies (Volume 1: Long Papers)}, 2024, 5364–5376.
\bibitem{zaratiana2025gliner2}
Zaratiana U, Pasternak G, Boyd O et al. GLiNER2: Schema-Driven Multi-Task Learning for Structured Information Extraction. In: \emph{Proceedings of the 2025 Conference on Empirical Methods in Natural Language Processing: System Demonstrations}, 2025, 130–140.
\bibitem{grattafiori2024llama3}
Grattafiori A, Dubey A, Jauhri A et al. The Llama 3 Herd of Models. \emph{arXiv} 2024;arXiv:2407.21783.
\bibitem{jiang2023mistral}
Jiang AQ, Sablayrolles A, Mensch A et al. Mistral 7B. \emph{arXiv} 2023;arXiv:2310.06825.
\bibitem{abdin2024phi3}
Abdin M, Aneja J, Awadalla H et al. Phi-3 Technical Report: A Highly Capable Language Model Locally on Your Phone. \emph{arXiv} 2024;arXiv:2404.14219.
\bibitem{guo2025deepseekr1}
Guo D, Yang D, Zhang H et al. DeepSeek-R1 incentivizes reasoning in LLMs through reinforcement learning. \emph{Nature} 2025;645:633–638.
\bibitem{gemma2026gemma4}
Gemma Team, El Abd S, Aggarwal V et al. Gemma 4 Technical Report. \emph{arXiv} 2026;arXiv:2607.02770.
\bibitem{qwen35card}
Qwen, 2026-03-02 (accessed 2026-11-07): Qwen/Qwen3.5-4B model card, Hugging Face (repository), https://huggingface.co/Qwen/Qwen3.5-4B.
\bibitem{deepseek2026v4}
DeepSeek-AI, Xu A, Lin B et al. DeepSeek-V4: Towards Highly Efficient Million-Token Context Intelligence. \emph{arXiv} 2026;arXiv:2606.19348.
\bibitem{deepseekv4flashcard}
DeepSeek-AI, 2026-08-01 (accessed 2026-11-07): deepseek-ai/DeepSeek-V4-Flash-0731 model card, Hugging Face (repository), https://huggingface.co/deepseek-ai/DeepSeek-V4-Flash-0731.
\bibitem{deepseekv41card}
DeepSeek-AI, 2026-10-16 (accessed 2026-11-07): deepseek-ai/DeepSeek-V4.1-Flash model card, Hugging Face (repository), https://huggingface.co/deepseek-ai/DeepSeek-V4.1-Flash.
\bibitem{zheng2023mtbench}
Zheng L, Chiang W, Sheng Y et al. Judging LLM-as-a-Judge with MT-Bench and Chatbot Arena. \emph{arXiv} 2023;arXiv:2306.05685.
\bibitem{liu2019sbert}
Reimers N, Gurevych I. Sentence-BERT: Sentence Embeddings using Siamese BERT-Networks. In: \emph{Proceedings of the 2019 Conference on Empirical Methods in Natural Language Processing and the 9th International Joint Conference on Natural Language Processing (EMNLP-IJCNLP)}, 2019, 3980–3990.
\bibitem{minilmcard}
Sentence-Transformers, 2026-06-01 (accessed 2026-11-07): sentence-transformers/all-MiniLM-L6-v2 model card, Hugging Face (repository), https://huggingface.co/sentence-transformers/all-MiniLM-L6-v2.
\bibitem{kwon2023vllm}
Kwon W, Li Z, Zhuang S et al. Efficient Memory Management for Large Language Model Serving with PagedAttention. In: \emph{Proceedings of the 29th Symposium on Operating Systems Principles}, 2023, 611–626.
\bibitem{ollama2026}
Ollama, Release 0.32.9 (cluster version; accessed 2026-11-07): Ollama local large-language-model runtime, GitHub (repository), https://github.com/ollama/ollama (MIT license).
\bibitem{lin2023awq}
Lin J, Tang J, Tang H et al. AWQ: Activation-aware Weight Quantization for LLM Compression and Acceleration. \emph{arXiv} 2023;arXiv:2306.00978.
\bibitem{schoch2020taxonomy}
Schoch C, Ciufo S, Domrachev M et al. NCBI Taxonomy: a comprehensive update on curation, resources and tools. \emph{Database (Oxford)} 2020;2020:baaa062.
\bibitem{brown2015gene}
Brown G, Hem V, Katz K et al. Gene: a gene-centered information resource at NCBI. \emph{Nucleic Acids Res} 2015;43:D36–D42.
\bibitem{lipscomb2000mesh}
Lipscomb CE. Medical Subject Headings (MeSH). \emph{Bull Med Libr Assoc} 2000;88:265–6.
\bibitem{davis2026ctd}
Davis A, Wiegers T, Sciaky D et al. Comparative Toxicogenomics Database: update 2027. \emph{Nucleic Acids Res} 2026;gkag1041.
\bibitem{barrett2012bioproject}
Barrett T, Clark K, Gevorgyan R et al. BioProject and BioSample databases at NCBI: facilitating capture and organization of metadata. \emph{Nucleic Acids Res} 2012;40:D57–D63.
\bibitem{leinonen2011}
Leinonen R, Sugawara H, Shumway M et al. The Sequence Read Archive. \emph{Nucleic Acids Res} 2011;39:D19–D21.
\bibitem{vasilevsky2022mondo}
Vasilevsky N, Matentzoglu N, Toro S et al. Mondo: Unifying diseases for the world, by the world. \emph{medRxiv} 2022;2022.04.13.22273750 (preprint).
\bibitem{gargano2024hpo}
Gargano M, Matentzoglu N, Coleman B et al. The Human Phenotype Ontology in 2024: phenotypes around the world. \emph{Nucleic Acids Res} 2024;52:D1333–D1346.
\bibitem{bodenreider2004umls}
Bodenreider O. The Unified Medical Language System (UMLS): integrating biomedical terminology. \emph{Nucleic Acids Res} 2004;32:267D–270.
\bibitem{whetzel2011bioportal}
Whetzel P, Noy N, Shah N et al. BioPortal: enhanced functionality via new Web services from the National Center for Biomedical Ontology to access and use ontologies in software applications. \emph{Nucleic Acids Res} 2011;39:W541–W545.
\bibitem{cote2006ols}
Côté R, Jones P, Apweiler R et al. The Ontology Lookup Service, a lightweight cross-platform tool for controlled vocabulary queries. \emph{BMC Bioinformatics} 2006;7:97.
\bibitem{hripcsak2015ohdsi}
Hripcsak G, Duke JD, Shah NH et al. Observational Health Data Sciences and Informatics (OHDSI): Opportunities for Observational Researchers. In: \emph{MEDINFO 2015: eHealth-enabled Health}. Amsterdam: IOS Press, 2015, 574–79 (see note).
\bibitem{yilmaz2011mixs}
Yilmaz P, Kottmann R, Field D et al. Minimum information about a marker gene sequence (MIMARKS) and minimum information about any (x) sequence (MIxS) specifications. \emph{Nat Biotechnol} 2011;29:415–420.
\bibitem{janssens2018disbiome}
Janssens Y, Nielandt J, Bronselaer A et al. Disbiome database: linking the microbiome to disease. \emph{BMC Microbiol} 2018;18:50.
\bibitem{wilson1927}
Wilson E. Probable Inference, the Law of Succession, and Statistical Inference. \emph{J Am Stat Assoc} 1927;22:209–212.
\bibitem{efron1993bootstrap}
Efron B, Tibshirani R. An Introduction to the Bootstrap. Boston: Springer, 1993.
\bibitem{cohen1960kappa}
Cohen J. A Coefficient of Agreement for Nominal Scales. \emph{Educ Psychol Meas} 1960;20:37–46.
\end{thebibliography}

\end{document}
